%!TEX root = ../../main.tex
%% ===============================================================
%% 4.6 교전 유형 분류와 부가 정보
%% 파일: chapters/04_localization/06_fight_type.tex
%% 상위: chapters/04_localization/chapter.tex
%% 정의 라벨: sec:loc-classify
%% 참조 라벨: -
%% 포함 객체: -
%% 인용: -
%% 상태: 초안 (docs/OUTLINE.md 참고)
%% ===============================================================

\section{교전 유형 분류와 부가 정보}
\label{sec:loc-classify}

병합된 교전에는 두 축의 유형이 부여된다. \emph{규모}(fight scale)는 근접 쌍의 수로 정의하여 8 쌍 이상이면 한타(teamfight), 4 쌍 이상이면 소규모 교전(skirmish), 그 미만이면 픽(pick)으로 하고, \emph{맥락}(fight context)은 중심 위치가 바론·용·전령 둥지 1{,}500 단위 이내, 포탑 1{,}000 단위 이내, 본진 3{,}000 단위 이내인지에 따라 오브젝트 교전, 포탑 다이브, 본진 교전 등으로 나눈다. 또한 교전 중 반경 3{,}000 단위 안의 비-킬 상호작용(와드, 소환사 주문)과, 마지막 킬 이후 30 초 동안의 오브젝트·포탑·골드 변동을 부가 정보로 기록한다. 이들은 하위 집단 분석에만 쓰이며 라벨이나 입력이 아니다.
